\chapter{Chapter 2 Refresher: Optimization, Generalization, and Scaling Intuition --- Optimization and Scaling}

\section*{Setting the Scene}
Chapter~2 consumes the gradient of an empirical risk built in Chapter~1 (Preface Steps~5--8): we optimise an expectation estimated from batches. The novelty is size: billions of parameters force \emph{cheap, noisy} directional probes rather than exact curvature solves.

\paragraph{Step A --- Local linear model.}
\textbf{Problem.} Full loss landscapes are unwieldy; near $\theta_t$ we approximate $\mathcal{L}(\theta_t+\Delta)$ by first-order Taylor expansion.
\textbf{Insight.} Gradient descent assumes meaningful linear prediction within chosen step lengths.
\textbf{Lineage.} Cauchy (1847) gradient methods; modern proofs invoke convex-like intuitions rarely satisfied literally in deep nets.
\textbf{Mental model.} If curvature tightens, shrink steps or add momentum filters---otherwise linearisation lies.

\paragraph{Step B --- Stochastic gradients.}
\textbf{Problem.} Exact $\nabla\mathcal{L}$ over full corpora is unaffordable.
\textbf{Insight.} Subsample to inject noise whose variance trades memory for stability (Preface Step~6).
\textbf{Lineage.} Robbins--Monro (1951); Bottou's large-scale ML synthesis.
\textbf{Mental model.} Treat minibatch gradient as true gradient + zero-mean perturbation; diagnose spikes as breakdown of bounded-variance stories.

\paragraph{Step C --- Adaptive signal processing.}
\textbf{Problem.} Parameter blocks receive gradients at wildly different scales (embeddings vs logits vs biases).
\textbf{Insight.} Running moment estimates rescale updates elementwise while preserving descent intent at smooth regions.
\textbf{Lineage.} Polyak momentum; Duchi/Hazan/Singer (Adagrad); Kingma/Ba (Adam); Loshchilov/Hutter (AdamW).
\textbf{Mental model.} Adaptive methods behave like learned preconditioners---trust them until drift/outliers falsify stationarity.

\paragraph{Step D --- Numerical hygiene.}
\textbf{Problem.} Low-precision tensors amplify underflow/overflow.
\textbf{Insight.} Loss scaling preserves gradient magnitude while exploiting hardware throughput.
\textbf{Lineage.} NVIDIA mixed-precision training recipes (2017 era onward).
\textbf{Mental model.} Clip gradients when tail events violate FP16 comfort zones.

\paragraph{Step E --- Macroscopic empirical laws.}
\textbf{Problem.} Engineers must budget parameters \emph{and} tokens jointly.
\textbf{Insight.} Iso-loss contours summarize decades of runs---Preface Step~7 interpreted informally at massive $n$.
\textbf{Lineage.} Kaplan et al.\ (2020); Hoffmann et al.\ / Chinchilla (2022).
\textbf{Mental model.} Scaling fits extrapolate poorly across dataset contamination regimes---treat coefficients as local Taylor surrogates.

\paragraph{Connective thread.}
Exact gradients are expensive; minibatches trade bias for variance: write $\hat g_t=\nabla_\theta \mathcal{L}_t(\theta_t)+\varepsilon_t$ with $\mathbb{E}[\varepsilon_t]=0$. Under i.i.d.\ sampling, $\mathrm{Var}(\varepsilon_t)$ scales roughly like $1/B$ with batch size $B$, which is why larger batches smooth updates but cost memory. Momentum and adaptive moments are filters on that noise; clipping guards rare spikes that violate the ``bounded variance'' story.

\section*{Notation Introduced in This Chapter}
\begin{itemize}[leftmargin=1.5em]
  \item $\theta$: model parameters.
  \item $\mathcal{L}(\theta)$: objective (loss) as function of parameters.
  \item $\nabla\mathcal{L}(\theta)$: gradient vector.
  \item $\eta_t$: learning rate at step $t$.
  \item $\hat g_t$: stochastic (minibatch) gradient estimate.
  \item $m_t, v_t$: first and second moment EMAs.
  \item $\beta_1,\beta_2$: EMA decay coefficients.
  \item $\tau$: clipping threshold.
  \item $N,D$: model/data scale variables used in scaling-law fits.
\end{itemize}

\section*{What You Need To Recall Precisely}
\begin{itemize}[leftmargin=1.5em]
  \item \textbf{Gradient as first-order local direction.} If $\mathcal{L}(\theta)$ is differentiable, then for small $\Delta$:
  \[
    \mathcal{L}(\theta+\Delta)\approx \mathcal{L}(\theta) + \nabla \mathcal{L}(\theta)^\top \Delta.
  \]
  Choosing $\Delta=-\eta\nabla\mathcal{L}$ gives first-order descent.

  \item \textbf{Stochastic gradient descent (SGD):}
  \[
    \theta_{t+1}=\theta_t-\eta_t \nabla \mathcal{L}(\theta_t).
  \]
  In practice, $\nabla \mathcal{L}$ is replaced by minibatch estimate $\hat g_t$.

  \item \textbf{Momentum as exponential averaging.}  
  With $m_t=\beta m_{t-1} + (1-\beta)\hat g_t$, updates use smoothed direction $m_t$.
  Think of $m_t$ as a low-pass filter on the gradient signal: high-frequency minibatch noise is attenuated while consistent downhill curvature survives, similar in spirit to a heavy ball continuing through shallow oscillations in the loss surface.

  \item \textbf{Bias correction for startup EMA.}  
  Because $m_0=0$, early EMAs are biased toward zero---the running average ``forgets'' that gradients were nonzero before time zero. Adam uses debiased estimates
  \[
    \hat m_t=\frac{m_t}{1-\beta_1^t},\qquad \hat v_t=\frac{v_t}{1-\beta_2^t}.
  \]
  so that if $\hat g$ were constant, $\hat m_t$ would recover that constant after a few steps instead of ramping up spuriously slowly.

  \item \textbf{AdamW intuition (brief).}
  Classic Adam entangles weight decay with adaptive preconditioning in ways that distort the intended $\ell_2$ penalty. AdamW applies weight decay \emph{directly} to $\theta$ (decoupled from the adaptive gradient transform), while still using $\hat m_t,\hat v_t$ to precondition the loss gradient step.


  \item \textbf{Gradient clipping rule:}
  \[
    g \leftarrow g \cdot \min\left(1,\frac{\tau}{\|g\|}\right).
  \]
  This bounds update magnitude without changing direction when clipping is active.

  \item \textbf{Mixed precision with loss scaling.}  
  Forward/backward often run in bf16/fp16 for throughput; tiny gradients can flush to subnormals or zero in those formats. Multiply the scalar loss by a large factor $s$ before backward pass so chain-rule gradients inherit scale; divide accumulated gradients by $s$ before the optimizer applies updates (often fp32)---recovering usable magnitude without changing the mathematically correct descent direction.

  \item \textbf{Scaling-law interpretation.}  
  Empirical forms such as
  \[
    L(N,D)=L_\infty + aN^{-\alpha} + bD^{-\beta}
  \]
  summarise measured iso-loss contours across runs when architecture and data pipelines stay comparable. Use them as \emph{budget planners}: given fixed compute $C\approx 6ND$ tokens-parameter conventions from Kaplan/Chinchilla-style analyses, pick $(N,D)$ pairs predicted to sit near the elbow rather than extrapolate blindly across regimes where data quality or objectives shifted.
\end{itemize}

\section*{How These Results Build on Each Other}
\begin{enumerate}[leftmargin=1.5em]
  \item First-order approximation motivates gradient descent.
  \item Minibatch estimation introduces noise, requiring variance management.
  \item Momentum/adaptive moments stabilize noisy trajectories.
  \item Bias correction fixes startup distortion in adaptive moments.
  \item Clipping and precision controls handle numerical and tail-event instability.
  \item Scaling fits summarize many such runs for planning decisions.
\end{enumerate}

\section*{Reminder Glossary (Term by Term)}
\begin{itemize}[leftmargin=1.5em]
  \item \textbf{Gradient:} local rate-of-change vector for loss.
  \item \textbf{Learning rate:} step size in parameter space.
  \item \textbf{Minibatch noise:} estimation error from finite sample gradients.
  \item \textbf{Momentum/EMA:} smoothed running estimate of gradient direction.
  \item \textbf{Second moment estimate:} running estimate of squared gradients.
  \item \textbf{Bias correction:} startup debiasing of EMAs.
  \item \textbf{Gradient clipping:} norm-bound safeguard against spikes.
  \item \textbf{Warmup:} gradual LR increase to stabilize early dynamics.
  \item \textbf{Scaling law:} empirical relation among model size, data, compute, and loss.
  \item \textbf{AdamW:} Adam-style adaptive steps with weight decay applied directly to parameters (decoupled from gradient preconditioning).
\end{itemize}

\section*{Worked Example (Training-Dynamics Diagnosis)}
Suppose a run exhibits:
\begin{itemize}[leftmargin=1.5em]
  \item training loss oscillation in first 2k steps,
  \item occasional gradient norm spikes from 12 to 140,
  \item mixed-precision overflow warnings.
\end{itemize}

Use the refresher's reasoning chain:
\begin{enumerate}[leftmargin=1.5em]
  \item \textbf{Diagnose instability source.} Early oscillation plus spikes suggests oversized effective step on high-curvature/noisy minibatches.
  \item \textbf{Apply clipping policy.} With threshold $\tau=20$, a spike at norm 140 is rescaled by factor $20/140\approx0.143$.
  \item \textbf{Check EMA startup behavior.} For $\beta_1=0.9$ and constant scalar gradient 2:
\[
m_1=0.1\cdot2=0.2,\quad m_2=0.9(0.2)+0.1(2)=0.38,\quad m_3=0.542.
\]
  Bias-corrected:
\[
\hat m_3=\frac{m_3}{1-0.9^3}=\frac{0.542}{0.271}\approx2.00.
\]
  \item \textbf{Stabilize numerics.} Increase loss-scaling headroom or use dynamic scaling rollback on overflow.
  \item \textbf{Interpret outcome.} If oscillation dampens while validation trend improves, intervention matched failure mechanism.
\end{enumerate}

\section*{Intuition Checkpoints}
\begin{itemize}[leftmargin=1.5em]
  \item Adam corrections are not magic; under drifting distributions they are approximate.
  \item Clipping controls step magnitude, not objective geometry.
  \item Mixed precision saves memory/throughput but shifts numerical failure modes.
  \item Scaling-law extrapolations degrade when data quality and objective shift.
\end{itemize}

\section*{Pointers}
\begin{itemize}[leftmargin=1.5em]
  \item Loss and KL notation: see Chapter 1 refresher.
  \item Matrix and norm geometry used by clipping and curvature intuition: see Chapter 4 refresher.
\end{itemize}

\section*{References and What To Use Them For}
\begin{itemize}[leftmargin=1.5em]
  \item \textbf{Goodfellow--Bengio--Courville, \emph{Deep Learning}} --- SGD, momentum, adaptive optimization intuition.
  \item \textbf{Kingma and Ba (Adam)} --- original adaptive-moment method and bias correction rationale.
  \item \textbf{CS229/CS231n notes} --- optimization diagnostics and practical training heuristics.
  \item \textbf{Hoffmann et al. (2022)} --- compute-optimal scaling interpretation and planning.
  \item \textbf{Bottou et al., optimization for large-scale ML} --- stochastic optimization theory background.
\end{itemize}
